% Chapter 9: what the thesis will leave behind.
% Keep the heading with its first lines rather than alone at a page foot.
\needspace{10\baselineskip}
\chapter{Expected contributions}\label{ch:contrib}

If the plan succeeds, the thesis will contribute the following.

\begin{enumerate}[label=C\arabic*, font=\sffamily\color{accent}, leftmargin=2.2em,
  itemsep=.5em]
\item \textbf{A method.} An exposure-weighted acquisition rule for
  spatial sensing, with its near-optimality guarantee and an open
  implementation that a city climate office can run on a laptop.
\item \textbf{Evidence on transfer.} The first systematic test of
  whether urban temperature models and placement policies carry over
  between cities, with guidance on which covariates to trust.
\item \textbf{A field result.} A measured, not simulated, estimate of
  what adaptive placement gains on real streets over a full summer,
  including the losses that simulations ignore.
\item \textbf{An open dataset.} Four months of street-level temperature
  and humidity from 60 relocating sensors, released with metadata under
  an open licence.
\end{enumerate}

Beyond these, the work gives cities a concrete answer to a question
they already face each spring: with a fixed budget, where should the
next sensor go so that the next heat warning reaches the people who
need it?
